\documentclass[12pt]{article}
\usepackage[utf8]{inputenc}
\usepackage[T1]{fontenc}
\usepackage{hyperref}
\title{Bruny Island Tasmanian Language Resources for K-12 Curriculum Development}
\author{Compiled by Claude Sonnet 5.5 (claude-sonnet-5-5) for Andr\'as Kornai}
\date{}
\begin{document}
\maketitle

\section{Introduction}
Bruny Island Tasmanian (ISO 639-3 xpz) is an extinct language of south-eastern Tasmania. Wikipedia reports a list of about 986 words and names the Milligan vocabulary as a source. The last likely speaker was Truganini, who died in 1876. There are no native speakers, and I found no school curriculum built on the historical language. The related mainland variety Nuenonne has a thinner record, from the Péron and D'Entrecasteaux expedition vocabularies. The reconstructed language palawa kani is a separate, modern project.

\section{Primary Educational Resources}
I found no primers or textbooks for Bruny Island Tasmanian.

\section{Linguistic and Reference Materials}
\begin{itemize}
\item Bruny Island Tasmanian language (Wikipedia): \url{https://en.wikipedia.org/wiki/Bruny_Island_Tasmanian_language}
\item Nuenonne language (Wikipedia): \url{https://en.wikipedia.org/wiki/Nuenonne_language}
\item Aboriginal \&amp; Torres Strait Islander Data Archive record (appeared in the search; I did not open it, so I have not confirmed its contents): \url{https://atsida.edu.au/node/351}
\end{itemize}

\section{Religious and Community Resources}
\begin{itemize}
\item palawa kani (Wikipedia): the Tasmanian Aboriginal Centre's reconstructed language, built from Plomley's 1976 word lists and Crowley and Dixon's 1981 chapter. \url{https://wikipedia.com/wiki/Palawa_kani}
\item Bruny Island community page on an Aboriginal story (appeared in search; not opened): \url{https://brunyisland.org.au/our-island/rodneydillon-aboriginal-story/}
\end{itemize}

\section{Digital and Audio Resources}
I found no audio recordings or software specific to Bruny Island Tasmanian. Native-Land.ca has a territory map for Nuenonne country: \url{https://native-land.ca/maps/territories/nuenonne/}

\section{Curriculum Development Context}
The historical record is a word list with no speakers. Schools would most appropriately work with the Tasmanian Aboriginal community and with palawa kani, which has its own governance. I did not verify any palawa kani teaching materials in this search.
\end{document}